%% MASS/SMASH paper — revised draft incorporating peer critique
%% Target: Journal of Geophysical Research: Space Physics  /  Solar Physics
%% LaTeX standard article, AGU-compatible notation

\documentclass[12pt,letterpaper]{article}

\usepackage{amsmath,amssymb,amsthm}
\usepackage{graphicx}
\usepackage{booktabs}
\usepackage{hyperref}
\usepackage{natbib}
\usepackage{geometry}
\usepackage{setspace}
\usepackage{xcolor}
\usepackage{siunitx}
\usepackage{algorithm}
\usepackage{algpseudocode}
\usepackage{multirow}

\geometry{margin=1in}
\doublespacing

\newcommand{\Bvec}{\mathbf{B}}
\newcommand{\dBvec}{\Delta\mathbf{B}}
\newcommand{\Bhat}{\hat{\mathbf{B}}}
\newcommand{\nT}{\,\mathrm{nT}}
\newcommand{\s}{\,\mathrm{s}}
\newcommand{\MDL}{\mathrm{MDL}}
\newcommand{\RSS}{\mathrm{RSS}}
\newcommand{\PVI}{\mathrm{PVI}}
\newcommand{\FDF}{\mathrm{FDF}}

\title{\textbf{MASS/SMASH: A Seam-Aware, MDL-Regularized Detector \\
for Magnetic Current Sheets in Turbulent Solar Wind}}

\author{
  Mac Mayo\\
  \small{SeamAware Research, independent}\\
  \small{\texttt{macmarlins@gmail.com}}
}

\date{\today}

\begin{document}

\maketitle

\begin{abstract}
We introduce MASS/SMASH (Multi-seAm Seam-Sensitive Modeling / Symmetric
Minimum-description-length Antipodal Seam Hunting), an algorithm for detecting
magnetic discontinuities in solar wind time series that combines two
physics-motivated candidate detectors --- an antipodal correlation scanner that
scores positions where $B(t) \approx -B(t+\Delta t)$ (antisymmetry consistent
with a field polarity reversal) and a roughness discontinuity detector that
identifies abrupt changes in local signal variance --- with a Minimum
Description Length (MDL) model-selection step that penalizes detections at an
explicit information cost of $\approx \log_2 N$ bits per seam.
A detection is accepted only when the compression gain from modelling a
structural break exceeds this cost, providing a principled mechanism for
conservative event selection.

We validate MASS/SMASH against 120 known injected sector-boundary crossings in
30 days of synthetic solar wind data (3-s cadence; 6 nT background field; 30\%
Kolmogorov turbulence), benchmarked against a partial-variance-of-increments
(PVI) detector and a $|d\mathbf{B}/dt|$ threshold baseline.
Against the injected crossings, MASS/SMASH achieves a precision of 0.921
(false discovery fraction 7.9\%) and recall of 0.483 (F1 = 0.634), compared to
PVI's precision 0.989 / recall 0.750 (F1 = 0.853) and the baseline's
precision 0.909 / recall 1.000 (F1 = 0.952).
MASS/SMASH is the most conservative of the three methods, detecting 63 events
with only 5 false detections; the baseline produces 132 detections with 12
false positives.

Magnetic field rotation-angle analysis reveals that MASS/SMASH detections are
systematically offset $\approx 30\s$ from the true polarity center: the algorithm
locates the onset of the tanh transition profile (steepest gradient) rather than
the sign-flip midpoint.
When rotation angles are computed at the matched injected crossing positions
rather than at the detection positions, MASS/SMASH true-positive events have
mean angles of $127^\circ$--$134^\circ$ at 30--120 s averaging windows,
statistically indistinguishable from both the PVI catalog (130°--134°) and
the full injected set (128°--133°).
The previously apparent depletion of near-$\pi$ events in MASS/SMASH detections
is therefore an artifact of detection localization, not evidence of different
physical event classes.
We discuss the implications for multi-method current-sheet surveys and outline
extensions to full 3-component detection and real spacecraft data.
\end{abstract}

\textbf{Keywords:} solar wind turbulence; magnetic current sheets; magnetic
reconnection; heliospheric current sheet; Minimum Description Length;
information theory; change-point detection; partial variance of increments

\clearpage

%% ============================================================
\section{Introduction}
\label{sec:intro}
%% ============================================================

Localized magnetic field discontinuities --- here called \emph{current sheets}
in the broad sense of coherent structures with strong spatial field gradients and
substantial directional rotation --- are among the most energetically important
features of the turbulent solar wind.
They concentrate free magnetic energy, serve as preferential sites of magnetic
reconnection and particle heating, and produce the intermittent signature of the
turbulent energy cascade at MHD and kinetic scales
\citep{Matthaeus1986,Servidio2009,Retino2007,Osman2014,Chasapis2015}.
Their occurrence rate ($\sim$1--4 per hour at 1 AU), thickness
($\sim 10^2$--$10^3$ km), and statistical relationship to plasma heating have
been extensively studied with Wind, ACE, Cluster, MMS, and Parker Solar Probe
data \citep{Greco2008,Greco2009,Osman2011,Perri2012,Chasapis2015}.

In practice, automated detection of current-sheet crossings relies on a small
number of algorithmic approaches.
The Partial Variance of Increments \citep[PVI;][]{Greco2008} forms the
dimensionless ratio
\begin{equation}
  \PVI(t) \;=\; \frac{|\dBvec(t,\tau)|}{\sqrt{\langle |\dBvec|^2 \rangle_T}},
  \label{eq:pvi}
\end{equation}
where $\dBvec(t,\tau) = \Bvec(t+\tau) - \Bvec(t)$ is the lagged increment and
the denominator is a running-window variance normalization.
PVI peaks above a threshold (typically $3\sigma$) are used to identify
coherent structures including current sheets and flux tube boundaries
\citep{Greco2008,Osman2011,Osman2014}.
Simpler threshold methods on $|d\mathbf{B}/dt|$ or $|\mathbf{B}|$ are also
widely used as rapid screeners.
Wavelet-based and spectral methods have been applied to characterize
intermittency without constructing an explicit event catalog
\citep{Bruno2013}.

These methods share a structural feature: detection is driven by the magnitude
of a local gradient or fluctuation.
A threshold on $\PVI$ or $|d\mathbf{B}/dt|$ will fire wherever the field
changes rapidly, regardless of whether the surrounding signal structure changes
character across the event.
In strongly turbulent intervals, gradient-based methods can produce substantial
false discovery fractions; in our benchmark, the $|d\mathbf{B}/dt|$ detector
produces 132 detections for 120 true crossings with 12 false detections ($\FDF
= 9.1\%$), and PVI produces 91 detections with 1 false detection ($\FDF = 1.1\%$)
but misses 25\% of the injected crossings.

We introduce an alternative motivated by the \emph{structural} signature of a
current sheet: the signal on opposite sides of the crossing is locally
antisymmetric.
MASS/SMASH detects crossings as \emph{seams} --- positions where the time series
satisfies $B(t) \approx -B(t+\Delta t)$ and where the signal variance changes
abruptly --- and accepts only detections where an explicit information-theoretic
cost criterion is satisfied.
Rather than thresholding a derived local statistic, the algorithm poses detection
as a model-selection problem: it compares the Minimum Description Length (MDL)
of fitting the signal as a single homogeneous process against the MDL of fitting
two segments separated by a structural break, with the seam penalized at
$\sim \log_2 N$ bits.
A seam is retained only when the description-length reduction exceeds this cost.

This formulation is related to a broader class of MDL/BIC changepoint methods
\citep{Rissanen1978,Schwarz1978} and piecewise-model structural-break detection
methods that are well established in time-series analysis.
The specific novelty of MASS/SMASH, to the best of our knowledge, is the
combination of an antisymmetry (antipodal correlation) detector with a roughness
discontinuity detector and a model-zoo MDL selection step, applied to magnetic
current-sheet identification in solar wind data.
We make no claim that this is the first application of MDL to time-series
changepoint detection; we do claim it is a new application of the seam/antipodal
framework to solar wind plasma physics.

The paper is organized as follows.
Section~\ref{sec:algorithm} describes the MASS/SMASH algorithm.
Section~\ref{sec:validation} describes the synthetic validation data, evaluation
against known injected crossings, and the comparison detectors.
Section~\ref{sec:results} presents detection performance and the rotation angle
analysis.
Section~\ref{sec:discussion} interprets the localization finding, discusses
complementarity with PVI, and outlines limitations.
Section~\ref{sec:conclusions} summarizes the main contributions.

%% ============================================================
\section{The MASS/SMASH Algorithm}
\label{sec:algorithm}
%% ============================================================

MASS/SMASH processes a univariate time series $y_1, \ldots, y_N$ and returns
a set of seam locations $\{\tau_1, \ldots, \tau_m\}$ with associated MDL
gain scores.
The pipeline has four stages: (1) candidate proposal from two detectors,
(2) non-maximum suppression (NMS), (3) per-configuration MDL scoring, and
(4) global MDL selection.

\subsection{Antipodal Correlation Detector}
\label{subsec:antipodal}

The antipodal correlation detector scores each candidate position $\tau$ by:
\begin{equation}
  \mathcal{A}(\tau) \;=\; \mathrm{corr}\!\left(y[\tau-w:\tau],\;
    -\,y[\tau:\tau+w]\right),
  \label{eq:antipodal}
\end{equation}
where $w$ is the half-window size (default 20 samples) and both windows are
$z$-score normalized before computing the Pearson correlation.
High $\mathcal{A}(\tau)$ indicates the signal before $\tau$ is strongly
anticorrelated with the signal after $\tau$ --- the defining antisymmetry of
a field polarity reversal.

For an idealized tanh-profile crossing centred at $\tau_0$,
$B(t) = B_0 \tanh((t-\tau_0)/\delta)$, the antipodal score is maximized where
the two half-windows symmetrically straddle the transition.
In the presence of turbulent noise, the score is reduced in proportion to the
signal-to-noise ratio of the reversal amplitude over the turbulence variance.
Candidates are local maxima of $\mathcal{A}$ above a threshold
(default 0.35), retained after NMS with minimum separation 60 s.

\subsection{Roughness Discontinuity Detector}
\label{subsec:roughness}

The roughness at position $i$ is $\mathcal{R}(i) = \sigma(\Delta y[i:i+w_r])$,
where $\Delta y$ denotes first differences and $w_r$ is the roughness window
(default 20 samples).
The detector score is the normalized roughness change:
\begin{equation}
  \mathcal{D}(\tau) \;=\;
    \frac{|\mathcal{R}(\tau) - \mathcal{R}(\tau-1)|}{\sigma(\mathcal{R})}.
  \label{eq:roughness}
\end{equation}
Candidates are local maxima of $\mathcal{D}$ above threshold 0.20, again with
NMS.
The roughness detector is complementary to the antipodal scanner: it fires where
the local variance structure changes, which can include the edges of the
current-sheet transition layer even when the antipodal symmetry is diluted by
strong turbulence.

\subsection{MDL Model Selection}
\label{subsec:mdl}

After merging the top-$k^*$ candidates from both detectors (default $k^*=3$
per detector), MASS/SMASH scores all seam configurations (subsets of at most
$m_\mathrm{max}$ candidates) with the MDL objective in bits:
\begin{equation}
  \MDL(m, \hat{\boldsymbol{\theta}}) \;=\;
    \underbrace{\frac{N}{2}\log_2\frac{\RSS}{N}}_{\text{data fit}}
    \;+\;
    \underbrace{\frac{p}{2}\log_2 N}_{\text{model complexity}}
    \;+\;
    \underbrace{\frac{\alpha}{2}\,m\,\log_2 N}_{\text{seam encoding}},
  \label{eq:mdl}
\end{equation}
where $\RSS$ is the pooled residual sum of squares across segments, $p$ is the
total model parameter count, $m$ is the number of seams, and $\alpha = 2.0$
is the seam penalty coefficient.
The seam encoding term reflects the information cost of recording $m$ seam
positions from $N$ possible locations ($\approx m\log_2 N$ bits) and explicitly
encodes the principle that each seam must compress the data by more than the
cost of specifying its location.

Within each segment, MASS/SMASH evaluates a model zoo: Fourier series at
three complexity levels ($K = 4, 8, 12$ harmonics), polynomials of degree
2, 3, and 5, and autoregressive models of orders 5, 10, and 15.
The best model per segment is selected by its BIC contribution; the pooled
$\RSS$ and $p$ enter equation~\ref{eq:mdl}.
The configuration minimizing total MDL is returned; a no-seam configuration
($m=0$) is always included, so a detection requires an explicit MDL improvement.

\subsection{Chunked Inference on Long Time Series}
\label{subsec:chunked}

For 30 days at 3-s cadence ($N = 864{,}000$ samples), MASS/SMASH is applied
in a sliding-window fashion: 600-s windows (200 samples) at 300-s steps,
yielding 8,638 overlapping chunks.
Each chunk is processed independently with $m_\mathrm{max}=1$.
Each sample accumulates a cumulative MDL gain across all windows in which it
was selected as a seam.
A global NMS pass with 300-s minimum separation then retains peaks above
a gain threshold.
We use 20 bits for the main detection benchmark (requiring the seam to earn
its information cost by a wide margin) and 2 bits for exploratory analysis.

Algorithm~\ref{alg:chunked} summarizes the full pipeline.

\begin{algorithm}[t]
\caption{MASS/SMASH chunked inference}
\label{alg:chunked}
\begin{algorithmic}[1]
\Require Signal $y[0:N]$, window $W$, step $S$, gain threshold $G$
\State Initialize $\mathrm{gain}[0:N] \gets 0$
\For{$\mathrm{start} \in \{0, S, 2S, \ldots, N-W\}$}
  \State Run MASS/SMASH on $y[\mathrm{start}:\mathrm{start}+W]$ with $m_\mathrm{max}=1$
  \If{best solution has $m>0$ and MDL gain $g>0$}
    \For{each local seam index $\tau_\ell$}
      \State $\mathrm{gain}[\mathrm{start}+\tau_\ell] \mathrel{+}= g$
    \EndFor
  \EndIf
\EndFor
\State Retain positions with $\mathrm{gain} > G$; NMS with $\Delta\tau_\mathrm{min} = 300\s$
\State \Return detected positions and scores
\end{algorithmic}
\end{algorithm}

%% ============================================================
\section{Validation Methodology}
\label{sec:validation}
%% ============================================================

\subsection{Synthetic Solar Wind Data}
\label{subsec:synth}

We validate MASS/SMASH on a 30-day, 3-s-cadence synthetic magnetic field time
series ($N = 864{,}000$ samples) calibrated to mimic Wind/MFI observations during
the March 2001 interval studied by \citet{Osman2014}.

The background field follows a sector-boundary model.
The nominal field direction is
$\hat{\mathbf{b}}_0 = (0.201, 0.201, 0.958)$ (after normalization), giving a
field strength of $B_0 = 6\nT$ and placing $\approx 96\%$ of the polarity
reversal amplitude in the Bz component, which is the component input to MASS/SMASH
in the benchmark.
One hundred and twenty heliospheric current sheet (HCS) crossings are placed at
Poisson-random times with mean rate $\sim$4 per day and minimum spacing of one
hour, consistent with the observed rate near 1 AU
during moderately active intervals \citep{Wilcox1965}.
At each crossing the polarity transitions smoothly via
\begin{equation}
  \mathrm{pol}(t) \;=\; \tanh\!\left(\frac{t - t_\mathrm{cross}}{\delta/3}\right)
  \label{eq:tanh}
\end{equation}
with transition widths $\delta$ drawn uniformly from $[15\s, 60\s]$, comparable
to observed current sheet thicknesses at 1 AU \citep{Vasquez2007}.

Turbulence is superimposed with power spectral density $P(f) \propto |f|^{-5/3}$
(Kolmogorov spectrum), amplitude $\sigma_\mathrm{turb} = 0.30 B_0 = 1.8\nT$ per
component, consistent with typical amplitude ratios $\delta B/B \sim 0.2$--$0.4$
in the inertial range at 1 AU \citep{Bruno2013}.
The resulting field has $|\mathbf{B}|_\mathrm{rms} = 6.5\nT$,
$\sigma_{|\mathbf{B}|} = 1.6\nT$, and $\sigma_{B_z} = 6.0\nT$.

The 120 injected crossing centers are recorded at the exact tanh midpoints
(where $\mathrm{pol}(t) = 0$) and used as the primary ground truth in all
performance evaluations.

\subsection{Comparison Detectors}
\label{subsec:comparators}

We compare MASS/SMASH against two detectors.

\textbf{PVI detector.}
PVI (equation~\ref{eq:pvi}) is computed from the full 3-component field with
lag $\tau = 1$ sample (3 s) and a 90-s normalization window.
Crossings are identified as peaks with $\PVI > 3.0$ and minimum inter-event
separation of 300 s, yielding 91 detections for the 30-day interval.
This PVI threshold corresponds approximately to the 99th percentile of the
PVI distribution in our data and is the standard choice for strong
discontinuities in solar wind studies \citep{Greco2008,Osman2011}.
PVI is simultaneously a detector (evaluated against injected crossings) and
a reference catalog (for secondary comparisons); the 91 PVI events are not
used as the primary ground truth.

\textbf{$|d\mathbf{B}/dt|$ baseline.}
The baseline applies a threshold of $\mu + 2.5\sigma$ to the time derivative
of $|\mathbf{B}|$, with minimum inter-event separation of 300 s.
This represents the simplest operational screening method.

\subsection{Evaluation Metrics}
\label{subsec:metrics}

A detected event at position $\hat{\tau}$ is a true positive (TP) if it falls
within a tolerance window of a known injected crossing center; each crossing
is matched at most once.
We report precision $P = \mathrm{TP}/(\mathrm{TP}+\mathrm{FP})$,
recall $R = \mathrm{TP}/(\mathrm{TP}+\mathrm{FN})$, F1, and the
false discovery fraction $\FDF = 1 - P$.
Two tolerance levels are used: 150 s (50 samples) to accommodate the geometric
offset between the PVI gradient peak and the sign-flip center (Section~\ref{subsec:offset}),
and 30 s (10 samples) as a strict reference.

\subsection{Rotation Angle Analysis}
\label{subsec:rotation_method}

For each detection at sample index $p$, the magnetic field rotation angle is:
\begin{align}
  \Bhat_\mathrm{before} &=
    \frac{\langle\Bvec\rangle_{[p-w_a:p]}}{|\langle\Bvec\rangle_{[p-w_a:p]}|}, \quad
  \Bhat_\mathrm{after} \;=\;
    \frac{\langle\Bvec\rangle_{[p:p+w_a]}}{|\langle\Bvec\rangle_{[p:p+w_a]}|},
  \label{eq:Bhat} \\
  \omega &= \arccos\!\left(\mathrm{clip}(\Bhat_\mathrm{before}
    \cdot \Bhat_\mathrm{after}, -1, 1)\right),
  \label{eq:omega}
\end{align}
with averaging windows $w_a \in \{10, 30, 60, 120\}\s$.
We compute $\omega$ both at the raw detection position $p$ and, for matched
true positives, at the corresponding injected crossing center, in order to
separate localization effects from physical differences between detector
populations.

Statistical comparison of angle distributions between detectors uses a
two-sample Kolmogorov-Smirnov (KS) test.
Phase-randomized surrogates (100 realizations; independent Fourier phase
randomization per component) provide a null model for testing whether
distributional differences are explainable by spectral properties alone.
Because only 100 surrogates are computed, the minimum reportable surrogate
$p$-value is $p < 0.01$.

%% ============================================================
\section{Results}
\label{sec:results}
%% ============================================================

\subsection{Detection Performance Against Injected Crossings}
\label{subsec:detection}

Table~\ref{tab:main_results} summarizes performance for all three detectors
against the 120 injected crossing centers.

\begin{table}[t]
\centering
\caption{Detection performance against 120 known injected crossing centers.
FDF = false discovery fraction = $1 - P$.}
\label{tab:main_results}
\begin{tabular}{lcccccccc}
\toprule
Detector & Tol. & $N_\mathrm{det}$ & TP & FP & $P$ & $R$ & F1 & FDF \\
\midrule
MASS/SMASH (MDL gain $> 20$ bits) & 150 s & 63  & 58 & 5  & 0.921 & 0.483 & 0.634 & 7.9\% \\
MASS/SMASH (MDL gain $> 20$ bits) & 30 s  & 63  & 18 & 45 & 0.286 & 0.150 & 0.197 & 71.4\% \\
\midrule
PVI ($\PVI > 3.0$)                & 150 s & 91  & 90 & 1  & 0.989 & 0.750 & 0.853 & 1.1\%  \\
PVI ($\PVI > 3.0$)                & 30 s  & 91  & 90 & 1  & 0.989 & 0.750 & 0.853 & 1.1\%  \\
\midrule
$|d\mathbf{B}/dt|$ baseline       & 150 s & 132 & 120 & 12 & 0.909 & 1.000 & 0.952 & 9.1\% \\
$|d\mathbf{B}/dt|$ baseline       & 30 s  & 132 & 120 & 12 & 0.909 & 1.000 & 0.952 & 9.1\% \\
\bottomrule
\end{tabular}
\end{table}

Several findings are evident from Table~\ref{tab:main_results}.

\textbf{The $|d\mathbf{B}/dt|$ baseline achieves perfect recall.}
All 120 injected crossings are matched within 150 s; the 12 false detections
correspond to high-gradient turbulent fluctuations that did not originate from
injected crossings.
Because the baseline and the tanh crossings are both gradient-based, the
baseline performance is identical at the 30-s and 150-s tolerance levels.

\textbf{PVI achieves the highest precision but misses 25\% of injected crossings.}
PVI finds 90 of 120 injected crossings with only 1 spurious detection (a 1.1\%
false discovery fraction).
The 30 missed crossings are weak events --- those with smaller Bz amplitude at
the crossing center due to overlapping turbulent fluctuations --- that do not
exceed the $\PVI > 3.0$ threshold.
PVI performance is again identical at both tolerance levels because PVI peaks
closely coincide with the tanh inflection points.

\textbf{MASS/SMASH has the second-highest precision among the three methods.}
Of 63 MASS/SMASH detections, 58 (92.1\%) are within 150 s of an injected
crossing --- comparable in precision to both PVI and the baseline.
MASS/SMASH recall (0.483) is substantially lower than the baseline (1.000) or
PVI (0.750), reflecting the conservative MDL gain threshold (20 bits).
At strict 30-s tolerance, MASS/SMASH recall falls to 0.150, indicating that
most MASS/SMASH detections are offset from the crossing center by more than
30 s; Section~\ref{subsec:offset} explains this offset.

\textbf{MASS/SMASH finds crossings that PVI misses.}
Of the 63 MASS/SMASH detections matched to injected crossings, 24 correspond
to crossings not in the PVI catalog (i.e., crossings PVI missed).
Thus MASS/SMASH and PVI each capture a distinct subset of the injected population;
their joint catalog would cover 90 + 24 = 114 of 120 injected crossings.

\subsection{Detection Localization: A Systematic Offset from the Polarity Center}
\label{subsec:offset}

The fall to 0.150 recall at 30-s tolerance (versus 0.483 at 150 s) reveals that
most MASS/SMASH true positives are offset from the injected crossing center by
30--150 s.
To characterize this offset, we computed, for each true positive MASS/SMASH
detection, the displacement from the matched injected crossing center:
\begin{equation}
  \Delta\tau_i = \tau_i^\mathrm{MS} - \tau_i^\mathrm{inj},
\end{equation}
where $\tau_i^\mathrm{MS}$ is the MASS/SMASH detection position and
$\tau_i^\mathrm{inj}$ is the nearest injected crossing center within 150 s.

The offset distribution is approximately uniform across $[-50, +50]$ samples
($[-150\s, +150\s]$) but with notable clustering near $\Delta\tau \approx \pm 10$
samples ($\pm 30\s$).
This is consistent with a known geometric effect: the antipodal correlation
$\mathcal{A}(\tau)$ (equation~\ref{eq:antipodal}) is maximized where the signal
is most antisymmetric, which for the tanh profile
$B(t) \propto \tanh((t-t_0)/\delta)$ is not at the midpoint $t_0$ but rather
near the inflection points $t_0 \pm \delta/3$, where the signal
transitions between its asymptotic values and the symmetric plateau.
For transition widths $\delta \in [15\s, 60\s]$, the inflection points are
10--20 s from the midpoint, consistent with the observed clustering.

MASS/SMASH thus identifies the \emph{onset} of the transition profile --- where
the magnetic field begins its directional change --- rather than the polarity
midpoint.
This is a systematic, predictable offset that could be corrected in downstream
analysis by applying a transition-model-based centroid correction.
Alternatively, for applications requiring precise polarity-center localization,
a post-detection sharpening step using the sign-flip midpoint of the detected
segment would be appropriate.

\subsection{Rotation Angle Analysis}
\label{subsec:angles}

Table~\ref{tab:angles} compares the magnetic field rotation angles computed
at two positions: at the raw detection position ($\omega^\mathrm{detect}$),
and at the matched injected crossing center ($\omega^\mathrm{inj}$, for
MASS/SMASH true positives only).

\begin{table}[t]
\centering
\caption{Magnetic field rotation angle analysis. Columns:
$\bar{\omega}^\mathrm{detect}$ = mean angle at detection position;
$\bar{\omega}^\mathrm{inj}$ = mean angle at matched injected crossing center
(true positives only, $n=58$ for MASS/SMASH);
$f_\pi$ = fraction with $\omega > 160^\circ$;
KS $p$ = two-sample KS test of MASS/SMASH ($\omega^\mathrm{detect}$) vs.\ PVI.}
\label{tab:angles}
\begin{tabular}{lccccccc}
\toprule
 & \multicolumn{3}{c}{At detection position} &
   \multicolumn{2}{c}{At inj.\ crossing (TP only)} & \\
\cmidrule(lr){2-4}\cmidrule(lr){5-6}
Window & PVI $\bar{\omega}$ & MS $\bar{\omega}^\mathrm{detect}$ &
  KS $p$ & MS $\bar{\omega}^\mathrm{inj}$ & Inj.\ all $\bar{\omega}$ \\
\midrule
10 s  & 114° & 0.5° & $1.7\times10^{-44}$ & 101° & 108° \\
30 s  & 130° & 3.9° & $1.1\times10^{-42}$ & 128° & 128° \\
60 s  & 132° & 42°  & $4.1\times10^{-31}$ & 132° & 132° \\
120 s & 134° & 83°  & $2.1\times10^{-13}$ & 134° & 133° \\
\bottomrule
\end{tabular}
\end{table}

\textbf{At detection positions, MASS/SMASH events show anomalously small rotation angles.}
At 10-s and 30-s averaging windows, the mean angle at MASS/SMASH detection
positions is near $0^\circ$.
This anomaly is a direct consequence of the detection offset described in
Section~\ref{subsec:offset}: if the detection fires $\approx 10$ samples (30 s)
before the crossing center, then a 10-s (3-sample) averaging window on each
side of the detection position will straddle the transition profile rather than
bracket it.
The upstream window then averages a mixture of pre-transition and transition
field values, driving the mean field toward zero and producing a spurious small
rotation angle.
The KS statistics ($p < 10^{-13}$ at all window sizes) confirm that detection
positions show statistically distinct angle distributions compared to PVI events,
but this distinction is a localization artifact, not evidence of different
physical event classes.

\textbf{At matched injected crossing positions, MASS/SMASH events have angles
consistent with the full injected population and with PVI events.}
When $\omega$ is computed at the matched injected crossing center for the 58
MASS/SMASH true positives, the mean angles are $101^\circ$ (10 s window),
$128^\circ$ (30 s), $132^\circ$ (60 s), and $134^\circ$ (120 s) ---
statistically indistinguishable from both the PVI catalog values
($114^\circ$, $130^\circ$, $132^\circ$, $134^\circ$) and the full injected
population ($108^\circ$, $128^\circ$, $132^\circ$, $133^\circ$) at 30--120 s
windows.
The previously observed depletion of near-$\pi$ rotation events in MASS/SMASH
detections therefore reflects the detection position offset, not a physically
distinct sub-population of current sheet events.

\textbf{Prediction 1 verdict revised.}
An earlier analysis (prior to accounting for the detection offset) classified
the rotation angle result as AMBIGUOUS: the distribution of $\omega$ at
detection positions was highly significantly different from PVI (KS $p <
10^{-26}$), but not in the direction of a $\pi$-excess as hypothesized.
Corrected for the offset, the MASS/SMASH true-positive events have the same
angle distribution as PVI events; the apparent anomaly was a localization
artifact.
The verdict for the topological $\pi$-excess prediction is therefore
\textbf{NULL}: there is no evidence that the MDL selection criterion
preferentially selects large-rotation events over moderate-rotation events
among genuine current sheet crossings.
This is an informative null result that constrains the framework.

%% ============================================================
\section{Discussion}
\label{sec:discussion}
%% ============================================================

\subsection{Detection Characteristics of MASS/SMASH}
\label{subsec:char}

MASS/SMASH occupies a distinct operating point in the precision-recall space.
Its principal advantage over the $|d\mathbf{B}/dt|$ baseline is precision:
5 false detections versus 12 for the baseline (at 150-s tolerance), for comparable
false discovery fractions (7.9\% vs.\ 9.1\%).
Its principal disadvantage relative to the baseline is recall: 0.483 versus
1.000.
Compared to PVI, MASS/SMASH has slightly lower precision (0.921 vs.\ 0.989)
but achieves an important complementary benefit: it finds 24 crossings that PVI
misses.

The physical reason MASS/SMASH misses more crossings than the baseline is the
stringent MDL gain criterion (20 bits).
For the algorithm to accept a seam, the two-segment model must compress the
data substantially better than a single-segment model.
In strongly turbulent intervals, many current-sheet crossings are embedded in
high-variance background fluctuations that a single polynomial or AR model
cannot fit well regardless of seam placement, reducing the relative gain from
introducing the seam.
Conversely, for the 58 detected true positives, the transition is sufficiently
prominent relative to the local turbulence that the MDL criterion fires cleanly.

\subsection{The Detection Offset as a Physical Diagnostic}
\label{subsec:offset_interp}

The systematic $\approx 30\s$ offset between MASS/SMASH detections and the
true polarity center is a reproducible feature of the antipodal correlation
mechanism.
Rather than treating this as a deficiency, it can be leveraged as a diagnostic:
the detection position marks the \emph{transition onset} (where the field begins
rotating from its asymptotic value), while the true polarity center lies
$\approx 30\s$ downstream.
In principle, MASS/SMASH could be adapted to report the transition onset,
onset-to-center delay, and implied transition width as outputs, providing
morphological characterization of each crossing without the need for an
independent profile-fitting step.

\subsection{Complementarity with PVI}
\label{subsec:comp}

PVI and MASS/SMASH together would detect 114 of 120 injected crossings --- 95\%
of the full ground-truth population --- with a combined false discovery fraction
of approximately 6\%.
The 6 missed crossings are those for which neither method fires: the crossing
amplitude is too small for PVI, and the turbulence context is too noisy for
the MDL gain to be significant.

For routine survey work, the $|d\mathbf{B}/dt|$ baseline achieves perfect recall
at the cost of 9.1\% false detections, which is often acceptable.
MASS/SMASH is most valuable as a \emph{follow-on selector}: applied to the
$|d\mathbf{B}/dt|$ initial catalog, it could reduce false detections while
also flagging crossings that the baseline missed entirely (those in the 24-event
PVI-MASS/SMASH intersection that are absent from the baseline's false-positive list).
A practical multi-stage pipeline would be:
\begin{enumerate}
  \item Apply the baseline to obtain a high-recall ($>95\%$) initial catalog.
  \item Apply MASS/SMASH with a moderate gain threshold (e.g., 5 bits) to
    re-score each event and flag those with strong structural evidence.
  \item Use the MDL gain score as a continuous quality metric rather than a
    binary threshold.
\end{enumerate}

\subsection{On the Topological Prediction and the Null Result}
\label{subsec:nullresult}

The rotation angle analysis, once corrected for the detection offset, provides
a null result for Prediction 1: the MDL selection criterion does not
preferentially select current sheets with near-$\pi$ rotation angles.
This is an informative constraint.
The MDL criterion is sensitive to the compressibility advantage of a two-segment
model over a one-segment model, which depends on the amplitude of the structural
change relative to the noise, not specifically on the rotation angle.
A crossing with a 120° rotation at high SNR will score higher MDL gain than a
crossing with a 170° rotation at low SNR.
The null result therefore does not imply the topological motivation for the
antipodal correlator is incorrect; it implies that the MDL gain metric, as
implemented, does not directly proxy the rotation angle.

A direct topological test would require a version of the algorithm that
explicitly penalizes the seam acceptance on the basis of the 3D rotation angle
rather than purely on the 1D Bz MDL gain.
That extension is left for future work.

\subsection{Limitations}
\label{subsec:limits}

\textbf{Synthetic data only.}
All quantitative results are from synthetic data calibrated to match statistical
properties of Wind/MFI observations.
Real solar wind current sheets exhibit a wider range of transition profiles,
background field orientations, plasma $\beta$ conditions, and co-rotating
structures that are not represented in our single-regime model.
Application to Wind, ACE, Parker Solar Probe, and Solar Orbiter data is required
before quantitative performance claims can be extended to real observations.

\textbf{Single-component detection.}
MASS/SMASH processes only the Bz component.
Oblique crossings, where the polarity reversal is distributed across components,
will have lower SNR in Bz and therefore lower MDL gain.
Full 3-component MDL detection --- processing the full $\mathbf{B}$ vector or
$|\mathbf{B}|$ alongside Bz --- would improve recall for oblique events.

\textbf{PVI as a reference catalog.}
Although we use injected crossings as ground truth, the PVI catalog is also
reported as a secondary reference.
Readers comparing MASS/SMASH against PVI events rather than against injected
crossings should note that the metrics reflect ``agreement with PVI,'' not
detection accuracy.
MASS/SMASH achieves only F1 = 0.44 by this secondary metric, largely because
it detects 24 events that PVI misses but which are genuine injected crossings.

\textbf{Parameter calibration.}
The MDL gain threshold (20 bits), seam penalty ($\alpha = 2.0$), antipodal
threshold (0.35), and roughness threshold (0.20) are set heuristically.
A systematic grid search or Bayesian calibration against labeled training data
would provide a better-characterized precision-recall operating point and enable
quantitative comparison with other algorithms.

%% ============================================================
\section{Conclusions}
\label{sec:conclusions}
%% ============================================================

We have introduced MASS/SMASH, an MDL-regularized algorithm for detecting
magnetic current sheets as structural seams in univariate solar wind time series.
The algorithm combines an antipodal correlation detector, a roughness
discontinuity detector, and an information-theoretic model-selection step that
penalizes seams at $\sim \log_2 N$ bits each.
Validation against 120 known injected crossings in synthetic solar wind data
yields the following principal findings.

\begin{enumerate}

\item \textbf{MASS/SMASH achieves 92\% precision (FDF = 7.9\%) at 150-s
tolerance}, comparable to the $|d\mathbf{B}/dt|$ baseline (FDF = 9.1\%) and
PVI (FDF = 1.1\%), while detecting 63 events versus 132 for the baseline
and 91 for PVI.
Its recall of 48\% reflects the conservative 20-bit MDL gain criterion.

\item \textbf{MASS/SMASH is complementary to PVI.}
Of the 58 true-positive MASS/SMASH detections, 24 correspond to crossings not
in the PVI catalog.
A joint MASS/SMASH + PVI catalog would detect 95\% of injected crossings with
$\sim$6\% false discovery fraction.

\item \textbf{MASS/SMASH detections are systematically offset $\approx 30\s$
from the polarity center.}
The antipodal correlator fires at the transition onset (inflection point of the
tanh profile) rather than the sign-flip midpoint.
This is a predictable, correctable localization bias that could be turned into
a morphological diagnostic.

\item \textbf{When corrected for the detection offset, MASS/SMASH true-positive
events have the same rotation angle distribution as PVI events and the full
injected population.}
The previously apparent small-angle anomaly at detection positions was a
localization artifact.
The corrected verdict for the topological $\pi$-excess prediction (Prediction 1)
is NULL: the MDL gain criterion does not selectively prefer large-rotation events
over moderate-rotation events among genuine current sheet crossings.

\end{enumerate}

Future work will focus on: (a) validation on Wind/MFI, Parker Solar Probe
FIELDS, and Solar Orbiter MAG data; (b) extension to full 3-component or
vector-magnitude detection; (c) transition-onset to polarity-center correction
as a post-processing step; (d) systematic MDL gain threshold calibration; and
(e) cross-correlation of MASS/SMASH events with plasma reconnection signatures
(electron jets, Hall fields, enhanced particle heating) to characterize the
physical events selected by the antipodal/MDL criterion.

\section*{Acknowledgments}

The author thanks the Wind/MFI team \citep{Lepping1995} and the NASA/SPDF CDAWeb
service for public archive access used during algorithm development.
Data generation and analysis used NumPy \citep{Harris2020}, SciPy
\citep{Virtanen2020}, and Matplotlib \citep{Hunter2007}.

\section*{Data and Code Availability}

All code for MASS/SMASH, synthetic data generation, PVI catalog construction,
and benchmark analysis is available at
\url{https://github.com/macmayo1993/seam-aware-modeling}.
The benchmark reproduces fully from a single command using the included synthetic
data generator.

%% ============================================================
\bibliographystyle{agufull08}
\bibliography{references}

\begin{thebibliography}{99}

\bibitem[Bruno and Carbone(2013)]{Bruno2013}
Bruno, R., and V.~Carbone (2013),
The solar wind as a turbulence laboratory,
\textit{Living Rev. Solar Phys.}, \textbf{10}, 2,
doi:10.12942/lrsp-2013-2.

\bibitem[Chasapis et al.(2015)]{Chasapis2015}
Chasapis, A., et al. (2015),
Thin current sheets and associated electron heating in turbulent space plasma,
\textit{Astrophys.\ J.\ Lett.}, \textbf{804}, L1,
doi:10.1088/2041-8205/804/1/L1.

\bibitem[Gosling et al.(2005)]{Gosling2005}
Gosling, J.~T., et al. (2005),
Magnetic reconnection at the magnetopause: Stress balance and velocity
changes in the exhaust,
\textit{Geophys.\ Res.\ Lett.}, \textbf{32}, L03105,
doi:10.1029/2004GL021371.

\bibitem[Greco et al.(2008)]{Greco2008}
Greco, A., P.~Chuychai, W.~H.~Matthaeus, S.~Servidio, and P.~Dmitruk (2008),
Intermittent MHD structures and classical discontinuities,
\textit{Geophys.\ Res.\ Lett.}, \textbf{35}, L19111,
doi:10.1029/2008GL035454.

\bibitem[Greco et al.(2009)]{Greco2009}
Greco, A., W.~H.~Matthaeus, S.~Servidio, P.~Chuychai, and P.~Dmitruk (2009),
Statistical analysis of discontinuities in solar wind ACE data and
comparison with intermittent MHD turbulence,
\textit{Astrophys.\ J.\ Lett.}, \textbf{691}, L111,
doi:10.1088/0004-637X/691/2/L111.

\bibitem[Harris et al.(2020)]{Harris2020}
Harris, C.~R., et al. (2020),
Array programming with NumPy,
\textit{Nature}, \textbf{585}, 357--362,
doi:10.1038/s41586-020-2649-2.

\bibitem[Hunter(2007)]{Hunter2007}
Hunter, J.~D. (2007),
Matplotlib: A 2D graphics environment,
\textit{Comput.\ Sci.\ Eng.}, \textbf{9}(3), 90--95,
doi:10.1109/MCSE.2007.55.

\bibitem[Lepping et al.(1995)]{Lepping1995}
Lepping, R.~P., et al. (1995),
The Wind magnetic field investigation,
\textit{Space Sci.\ Rev.}, \textbf{71}, 207--229,
doi:10.1007/BF00751330.

\bibitem[Matthaeus and Lamkin(1986)]{Matthaeus1986}
Matthaeus, W.~H., and S.~L.~Lamkin (1986),
Turbulent magnetic reconnection,
\textit{Phys.\ Fluids}, \textbf{29}, 2513,
doi:10.1063/1.865487.

\bibitem[Osman et al.(2011)]{Osman2011}
Osman, K.~T., W.~H.~Matthaeus, A.~Greco, and S.~Servidio (2011),
Evidence for inhomogeneous heating in the solar wind,
\textit{Astrophys.\ J.\ Lett.}, \textbf{727}, L11,
doi:10.1088/2041-8205/727/1/L11.

\bibitem[Osman et al.(2014)]{Osman2014}
Osman, K.~T., et al. (2014),
Magnetic reconnection and intermittent turbulence in the solar wind,
\textit{Phys.\ Rev.\ Lett.}, \textbf{112}, 215002,
doi:10.1103/PhysRevLett.112.215002.

\bibitem[Perri et al.(2012)]{Perri2012}
Perri, S., et al. (2012),
Evidence of superdiffusive transport of electrons accelerated at
interplanetary shocks,
\textit{Astrophys.\ J.\ Lett.}, \textbf{750}, L2,
doi:10.1088/2041-8205/750/1/L2.

\bibitem[Phan et al.(2006)]{Phan2006}
Phan, T.~D., et al. (2006),
A magnetic reconnection X-line extending more than 390 Earth radii in the
solar wind,
\textit{Nature}, \textbf{439}, 175--178,
doi:10.1038/nature04393.

\bibitem[Retinò et al.(2007)]{Retino2007}
Retinò, A., et al. (2007),
In situ evidence of magnetic reconnection in turbulent plasma,
\textit{Nature Phys.}, \textbf{3}, 236--238,
doi:10.1038/nphys574.

\bibitem[Rissanen(1978)]{Rissanen1978}
Rissanen, J. (1978),
Modeling by shortest data description,
\textit{Automatica}, \textbf{14}, 465--471,
doi:10.1016/0005-1098(78)90005-5.

\bibitem[Schwarz(1978)]{Schwarz1978}
Schwarz, G.~E. (1978),
Estimating the dimension of a model,
\textit{Ann.\ Statist.}, \textbf{6}(2), 461--464,
doi:10.1214/aos/1176344136.

\bibitem[Servidio et al.(2009)]{Servidio2009}
Servidio, S., W.~H.~Matthaeus, M.~A.~Shay, P.~A.~Cassak, and P.~Dmitruk
(2009),
Magnetic reconnection in two-dimensional magnetohydrodynamic turbulence,
\textit{Phys.\ Rev.\ Lett.}, \textbf{102}, 115003,
doi:10.1103/PhysRevLett.102.115003.

\bibitem[Vasquez et al.(2007)]{Vasquez2007}
Vasquez, B.~J., C.~S.~Abramenko, D.~K.~Haggerty, and C.~W.~Smith (2007),
Electron-scale current sheets in the solar wind: Properties and
implications for particle acceleration,
\textit{J.~Geophys.\ Res.}, \textbf{112}, A11102,
doi:10.1029/2007JA012449.

\bibitem[Virtanen et al.(2020)]{Virtanen2020}
Virtanen, P., et al. (2020),
SciPy 1.0: Fundamental algorithms for scientific computing in Python,
\textit{Nature Methods}, \textbf{17}, 261--272,
doi:10.1038/s41592-020-0772-5.

\bibitem[Wilcox and Ness(1965)]{Wilcox1965}
Wilcox, J.~M., and N.~F.~Ness (1965),
Quasi-stationary corotating structure in the interplanetary medium,
\textit{J.~Geophys.\ Res.}, \textbf{70}, 5793--5805,
doi:10.1029/JZ070i023p05793.

\end{thebibliography}

\end{document}
